%!TEX root = ../../thesis.tex
\chapter{要素技術}
\label{chap:elemental}

本章では，本研究で用いる要素技術について述べる．

\section{メトリックマップに基づくナビゲーション}
\label{sec:metric_navigation}

% TODO: データ収集の教師として用いた，占有格子地図と自己位置推定に基づくナビゲーション（ROS navigation，move_base）を書く．

\section{視覚ナビゲーションの基盤モデル}
\label{sec:related_vnm}

視覚ナビゲーションでは，複数のロボットと環境から収集したデータを用いるViNT\cite{vint}や，Transformerと拡散モデルを用いて目標指向ナビゲーションと探索を統一的に扱う\NoMaD\cite{nomad}が提案されている．

ViNTは，現在と過去の観測画像およびゴール画像をEfficientNet\cite{efficientnet}で符号化し，Transformerで統合して，ゴールまでの時間的な距離と将来のwaypoint列を出力する\cite{vint}．
異なるロボットで収集した多様な走行データで学習することで，学習に用いていない環境やロボットにも適用できることが示されている．
ViNTでは，GPS waypointや左折・右折・直進といった高レベルの経路指示を入力とする適応も検討されているが，その指示は，事前に記録した経路の地図とその上での自己位置推定に基づいて算出される\cite{vint}．

\NoMaD は，ViNTと同様の観測エンコーダに，ゴール画像の入力を確率的に隠すゴールマスキングを導入したモデルである\cite{nomad}．
行動の生成には拡散モデル\cite{ddpm,diffusion_policy}を用い，ノイズから段階的にwaypoint列を生成することで，複数の選択肢がある場面での多峰的な行動分布を表現する．
\NoMaD はゴール画像をマスクすることでゴール非依存の探索も行えるが，その行動は特定の目的地へ向かうものではない．
本研究では，このゴールマスク時の観測表現に目標方向の情報を加えることで，ゴール画像を用いずに指定した方向へ進ませる．

\section{画像認識の基盤モデル}
\label{sec:related_vision}

画像認識においても，自己教師あり学習によって汎用的な視覚表現を獲得した\DINOvTwo \cite{dinov2}や，単眼RGB画像から絶対尺度の距離を推定する\UniDepth \cite{unidepth}が提案されている．

\DINOvTwo は，Vision Transformer（ViT）\cite{vit}を，ラベルを用いない自己蒸留によって大規模な画像データセットで学習したモデルである\cite{dinov2}．
得られた特徴は，特徴抽出器を固定したまま線形層を学習するだけで，画像分類やセグメンテーション，深度推定など多様なタスクに利用できることが示されている．
本研究では，この性質を利用し，\DINOvTwo を凍結したまま通路分類に用いる．

\UniDepth は，単眼RGB画像から，カメラの内部パラメータを与えずに絶対尺度の深度を推定するモデルである\cite{unidepth}．
カメラの情報を画像から推定するモジュールを持ち，多様なデータセットで学習することで，学習に用いていない環境でも距離を推定できる．
本研究では，\UniDepth の推定深度から，通路分類に用いる鳥瞰の占有格子と，衝突回避に用いる障害物の点群を求める．

\section{衝突回避}
\label{sec:related_care}

\CARE は，視覚ナビゲーションモデルが出力した軌道を，推定した障害物からの反発力に基づいて補正する衝突回避手法である\cite{care}．
ナビゲーションモデル自体を再学習せずに組み合わせられるため，本研究では経路追従モジュールの出力の補正に用いる．

\section{Pure Pursuit}
\label{sec:pure_pursuit}

% TODO: 経路追従モジュールで用いる Pure Pursuit（経路上の前方注視点へ向かう円弧の曲率から角速度を求める制御）を書く．
